% Generated draft: requires mathematical and source review.
\documentclass[11pt]{article}
\usepackage[margin=25mm]{geometry}
\usepackage{fontspec}
\setmainfont{Latin Modern Roman}
\usepackage{amsmath,amssymb,mathtools,unicode-math}
\setmathfont{Latin Modern Math}
\usepackage{enumitem,needspace,xurl,hyperref}
\setlength{\parindent}{0pt}
\setlength{\parskip}{5pt}
\setlength{\emergencystretch}{3em}
\newcommand{\N}{\mathbb{N}}
\newcommand{\Z}{\mathbb{Z}}
\newcommand{\Q}{\mathbb{Q}}
\newcommand{\R}{\mathbb{R}}
\newcommand{\C}{\mathbb{C}}
\newcommand{\F}{\mathbb{F}}
\newcommand{\A}{\mathbb{A}}
\newcommand{\PP}{\mathbb{P}}
\newcommand{\E}{\mathbb{E}}
\newcommand{\eps}{\varepsilon}
\newcommand{\dd}{\,\mathrm{d}}
\begin{document}
\section*{UnsolvedMath ID: 20002837}
\textbf{Source problem code:} catalogue entry
\subsection{Definitions and mathematical statement}
The source asks how moduli of two-dimensional representations relate to pseudorepresentations near a split residual representation $\omega_1\oplus\omega_2$ with distinct characters. It recalls a Nyssen--Rouquier equivalence in an absolutely irreducible setting, then sketches projective extension fibers $X_{\omega_i}$ over $\operatorname{Spec}R(T_{\mathbb F})$, possible gluing at the trivial extension, and pushforward of a universal rank-two bundle. The group, deformation conditions, precise moduli functors, and diagrams are not fully supplied. An equivalence of appropriate absolutely irreducible deformation functors must not be silently interpreted as an equivalence of unrigidified stacks with all scalar automorphisms.

The aid supplies a linear local diagnostic. Let
\[
 U=\operatorname{Ext}^1(\omega_1,\omega_2),\qquad
 W=\operatorname{Ext}^1(\omega_2,\omega_1)
\]
be finite-dimensional extension spaces in the intended representation category over a field $\mathbb F$. After removing the common scalar stabilizer, the effective diagonal stabilizer $\mathbb G_m$ acts with opposite weights:
\[
 t\cdot(u,w)=(tu,t^{-1}w),\qquad
 \mathfrak X=[(U\oplus W)/\mathbb G_m].
\]
If $\dim U=a$, $\dim W=b$, its affine invariant quotient is
\[
 \operatorname{Spec}\mathbb F[z_{ij}:1\le i\le a,1\le j\le b]/I_2(z),
 \qquad z_{ij}=u_iw_j,
\]
the rank-at-most-one determinantal cone. If one extension space is zero the quotient is a point. Mixed invariants are quadratic, with the two-by-two-minor relations.

The nonzero one-sided axes give
\[
 [(U\setminus\{0\})/\mathbb G_m]\cong\PP(U),\qquad
 [(W\setminus\{0\})/\mathbb G_m]\cong\PP(W),
\]
with empty projective space if the corresponding vector space is zero. Both map to the cone vertex and are proper over the field. Their zero extension is omitted, so it is not an ordinary scheme point at which these projective spaces can simply be glued. The split point in the effective quotient retains a residual gerbe $B\mathbb G_m$; an unrigidified representation stack also retains the removed common scalars. This explains the need for stack data without proving that this linear quotient is the entire actual formal deformation space: obstruction equations and diagonal-character deformations may be additional.

The original assertion that the proposed $X_{\omega_i}$ are proper needs their full definitions and a genuine valuative argument; properness of the model's projective axes alone does not establish it. Likewise proper pushforward does not automatically produce a vector bundle. A general dual-number obstruction is the proper closed immersion $i:\operatorname{Spec}\mathbb F\to\operatorname{Spec}\mathbb F[\varepsilon]/(\varepsilon^2)$: $\mathcal O^{\oplus2}$ is a rank-two bundle on the source, while $i_*\mathcal O^{\oplus2}$ is not locally free on the target. This illustrates the missing flatness/base-change requirements, not failure for a fully defined particular universal family. The aid's model and obstruction therefore diagnose the sketch but do not construct the requested global representation stack or its universal bundle.
\subsection{Short English statement}
The supplied model retains opposite extension weights and the split stabilizer, producing a rank-one cone and projective one-sided fibers. It explains why the proposed gluing is stack-theoretic and why proper pushforward alone is insufficient. It is a linear diagnostic, not an identification of the complete deformation moduli.
\subsection{Sources}
\begin{itemize}
\item Local UnsolvedMath registry, permanent ID 20002837, source code catalogue entry; created and updated 14 August 2026. Canonical \texttt{aim-representation-theory-notes.json}, zero-based \texttt{notes[108]}, tag \texttt{problem}, source item 2.
\item AIM ``$p$-adic representations, modularity, and beyond,'' section 8.2, Problem 2, according to the supplied verification: \url{https://aimath.org/WWN/padicmodularity/padicmodularity.pdf}. Statement status is reconstructed\_unverified. The verification corrects ``Rouquire'' to Rouquier; its final displayed-map recovery is truncated. Original remarks and literature field are empty.
\item The source attributes the absolutely irreducible theorem to Nyssen and Rouquier and records discussion by ``Mark'' and ``Richard,'' without full bibliographic data or identifying their surnames. Research attempt 1: \texttt{partial\_result}, problem status at run \texttt{partially\_solved}, rigor label \texttt{proved}, novelty confidence low.
\item Ulam AI / UnsolvedMath Contributors, CC BY 4.0. Machine-generated research aid, not peer reviewed; independently verify before citation. Complete authoritative references and actual deformation/obstruction hypotheses are absent.
\end{itemize}
\end{document}
